% Figure 4: Nested supercells under code truncation.
% Each panel shows the same 8x8 background grid; the blue rectangle is
% the supercell S(\hat c) for the code below it.
% From left to right: dropping a character zooms out (the supercell grows).
\begin{tikzpicture}[scale=0.42,>=Stealth,thick]

\foreach \i/\side/\xs/\label/\res in {%
  0/1/0/{\texttt{CR:4H7}}/{$\sim 19$\,mm},
  1/2/10/{\texttt{CR:4H}}/{$\sim 37$\,mm},
  2/4/20/{\texttt{CR:4}}/{$\sim 75$\,mm},
  3/8/30/{\texttt{CR}}/{$\sim 150$\,mm}}{%
  \begin{scope}[xshift=\xs cm]
    \draw[gray!40,step=1] (0,0) grid (8,8);
    \draw[blue!70,thick,fill=blue!15,fill opacity=0.5]
      (0,0) rectangle (\side,\side);
    \node[font=\normalsize\bfseries] at (4,-0.9) {\label};
    \node[font=\small,gray!50!black] at (4,-1.9) {\res};
  \end{scope}
}

% Arrows: dropping one character grows the supercell (zoom out).
\draw[->,thick,gray!70]
  (8.3,4) -- node[above,font=\footnotesize] {drop char} (9.7,4);
\draw[->,thick,gray!70]
  (18.3,4) -- node[above,font=\footnotesize] {drop char} (19.7,4);
\draw[->,thick,gray!70]
  (28.3,4) -- node[above,font=\footnotesize] {drop char} (29.7,4);

\end{tikzpicture}
